% ECONOMICS_PROBLEM: {"id": "H11-300", "title": "Origins and persistence of state capacity", "jel_code": "H11", "source_id": "OP-0127"}
% !TeX program = xelatex
\documentclass[11pt,a4paper]{article}
\usepackage[margin=23mm]{geometry}
\usepackage{fontspec}
\setmainfont{Latin Modern Roman}
\usepackage{amsmath,amssymb,mathtools}
\usepackage{xcolor,enumitem,booktabs,longtable,array,xurl,hyperref}
\definecolor{muted}{HTML}{53636C}
\hypersetup{colorlinks=true,urlcolor=blue,linkcolor=blue}
\setlength{\parindent}{0pt}
\setlength{\parskip}{5pt}
\newcommand{\R}{\mathbb R}
\newcommand{\E}{\mathbb E}
\newcommand{\N}{\mathbb N}
\newcommand{\ind}{\mathbf 1}
\newcommand{\Var}{\operatorname{Var}}
\newcommand{\Cov}{\operatorname{Cov}}
\newcommand{\supp}{\operatorname{supp}}
\DeclareMathOperator*{\argmin}{arg\,min}
\DeclareMathOperator*{\argmax}{arg\,max}
\newcommand{\entrymeta}[3]{{\sffamily\small\color{muted}\textbf{\texttt{#1}}\quad|\quad #2\par #3\par}}
\newcommand{\Problemref}[1]{\texttt{#1}}
\newcommand{\JELref}[2]{\texttt{#2} (JEL \texttt{#1})}
\begin{document}
\section*{Origins and persistence of state capacity}
\textbf{Problem ID:} \texttt{H11-300}\quad \textbf{JEL:} \texttt{H11}\par
% BEGIN ECONOMICS STATEMENT
\label{problem:OP-0127}
\label{atlas:prob:Q7}
{\sffamily\small\color{muted}\textbf{Q7} | Political economy | Editorial research problem family\par}
\subsection*{Definitions and assumptions}
Fix states $k_t\in\mathbb R_+^d$ describing measured fiscal, legal and administrative capabilities, a political state $p_t$, and a finite or discounted infinite horizon. Government investment $i_t\in\mathcal I(k_t,p_t)$ has cost $c(i_t,k_t)$ in a declared numeraire and transition
\[
k_{t+1}=F(k_t,i_t,p_t,\varepsilon_{t+1}).
\]
The transition preserves $\mathbb R_+^d$ and includes depreciation and complementarity if present. Citizen and elite decisions determine $p_{t+1}$ through a specified political game. A model $m\in\mathcal M$ defines shocks, private information, objective functions, feasibility, measurement and equilibrium selection. Require all discounted welfare sums to be finite. For every horizon used below, $\mathbb E_m[\|k_{t+h}(a)\|+|Y_{t+h}(a)|]<\infty$ under both compared interventions.
\subsection*{Formal research question}
Identify or bound the causal transition law under feasible investment policies $a$, the output and tax-capacity responses
\[
\theta_h(a,a_0)=\mathbb E_m[(k_{t+h}(a),Y_{t+h}(a))-(k_{t+h}(a_0),Y_{t+h}(a_0))],
\]
and the welfare-maximizing policy set. Determine which political incentives sustain high-capacity states, and whether historical and contemporary data distinguish complementarity-driven persistence from selection into capacity investment. If multiple political equilibria survive, give selector-specific or set-valued policy comparisons.
\subsection*{Interpretation and scope}
Capacity must be operationalized through service or revenue capabilities, not equated with realized tax revenue alone, which also depends on policy and the tax base. This question differs from D2's intervention design by emphasizing origins and persistence of the endogenous state. Historical war instruments require exclusion restrictions: war can affect development through destruction, borders, culture or selection as well as fiscal capacity. Fit of a dynamic model is not identification of those channels.
\subsection*{Source audit and status}
All three author-year citations are verified. The atlas's linked Besley--Persson source is the 2009 working paper corresponding to the 2010 Econometrica article; both dates are distinguished. Dincecco--Prado's journal is Journal of Economic Growth. These sources establish models and evidence under stated assumptions. The joint identification and policy target is editorial; no claim is made that their benchmark models remain unsolved.
\subsection*{References}
{\footnotesize\raggedright
\begin{enumerate}[label={\textbf{[S\arabic*]}},leftmargin=3em]
\item Timothy Besley and Torsten Persson (2009). \emph{The Origins of State Capacity: Property Rights, Taxation, and Politics}. American Economic Review 99(4), 1218--1244.\par
\textit{Locator and source role:} Abstract and investment model; legal and fiscal capacity complementarity under stated political assumptions.\par
\url{https://www.aeaweb.org/articles?id=10.1257/aer.99.4.1218}
\item Timothy Besley and Torsten Persson (2010). \emph{State Capacity, Conflict, and Development}. Econometrica 78(1), 1--34.\par
\textit{Locator and source role:} Abstract and capacity-investment model. Atlas link [25] is the 2009 NBER Working Paper 15088, not the 2010 publication.\par
\url{https://onlinelibrary.wiley.com/doi/10.3982/ECTA8073}
\item Mark Dincecco and Mauricio Prado (2012). \emph{Warfare, Fiscal Capacity, and Performance}. Journal of Economic Growth 17(3), 171--203.\par
\textit{Locator and source role:} Abstract and historical-war instrument; metadata and substantive question verified from the institutional research record.\par
\url{https://research.cbs.dk/en/publications/warfare-fiscal-capacity-and-performance/}
\end{enumerate}
}

\subsection*{Common conventions: Source mathematical conventions}

Unless an entry states otherwise, agent and firm sets are finite. State and action spaces are measurable, strategies and outcome maps are measurable, probability laws are specified, and expectations exist for the targets under discussion. Infinite-horizon discount factors lie in $(0,1)$ and discounted objectives are finite. Norms, tolerances and complexity input sizes must be specified for computational comparisons. Constants or primitives that appear locally are inputs, not empirical facts asserted by this catalogue.

\textbf{Identification.} A statistical model $m\in\mathcal M$ implies an observable law $P_m$ and target $\theta(m)$. At law $P$, its identified set is
\[
\Theta(P;\mathcal M)=\{\theta(m):m\in\mathcal M,\ P_m=P\}.
\]
Point identification holds when this set is a singleton, or equivalently when $P_{m_1}=P_{m_2}$ implies $\theta(m_1)=\theta(m_2)$ for all compatible models. Sharp bounds mean bounds attained, or approached when the boundary is not attained, by compatible models. Writing this set is a definition, not a solution: the substantive task is to establish informative restrictions and derive or estimate the set. An unrestricted-confounding model may give only trivial bounds.

\textbf{Inference.} Identification, estimability and valid uncertainty quantification are separate questions. Fix $\alpha\in(0,1)$. A confidence procedure $C_n$ over a declared law class $\mathcal P$ has asymptotic uniform coverage at level $1-\alpha$ when
\[
\liminf_{n\to\infty}\inf_{P\in\mathcal P}\Pr_P\{\theta(P)\in C_n\}\geq1-\alpha.
\]
For set-identified targets the coverage event must be defined separately, for example coverage of the full identified set. If an entry uses identification-indexed models instead of uniquely indexed laws, the same coverage convention is applied over those models. Pointwise limits do not establish uniform validity, and asymptotic validity does not guarantee finite-sample precision. Sample sizes, dependence structures and weak-identification sequences are part of each application.

\textbf{Counterfactuals.} $Y_i(a)$ is a potential outcome under a specified intervention, including its timing, implementation and versions. It is not $Y_i$ conditional on voluntarily choosing $a$. A full intervention vector permits interference; shorthand scalar treatment notation does not silently impose no interference. Assignment, selection, attrition, anticipation and measurement must be specified. A claim about an instrument's local effect is not automatically a population policy effect.

\textbf{Economic equilibrium.} A counterfactual model includes best responses, constraints, feasibility, market clearing, state transitions and expectations consistent with the stated information. When several equilibria exist, either a declared selector is an input or the target is set-valued. Comparisons across policy regimes must use the stated selection convention. Reduced-form relationships are not presumed invariant after an equilibrium-changing intervention.

\textbf{Optimization and welfare.} Optimization is over the stated feasible set. If attainment is not established, $\sup$ or $\inf$ and $\varepsilon$-optimal choices replace an asserted maximizer or minimizer. For example, with value $V=\sup_{a\in\mathcal A}W(a)<\infty$, an $\varepsilon$-optimal set is $\{a\in\mathcal A:W(a)\geq V-\varepsilon\}$ for $\varepsilon>0$. Benefits, costs, risk and health or environmental outcomes can be aggregated only using declared units and conversion weights. Interpersonal utility weights, the treatment of fiscal transfers and foreign profits, discounting, and the behavioral welfare criterion are explicit inputs. A comparative-static derivative additionally requires existence and appropriate differentiability of the selected counterfactual.

\textbf{Sources.} Local labels [S1], [S2], and so forth restart in every question. Locators identify the relevant abstract, model, section or result at the level actually checked; a bibliographic or abstract-level verification is not represented as inspection of an entire proof. Working papers are distinguished from later publications and changing versions. Source summaries are paraphrased. All URL checks and editorial formulations refer to the audit date stated above.
% END ECONOMICS STATEMENT
\end{document}
